% !TEX root = ../main.tex
\chapter{Introduction and Research Problem}
\label{ch:introduction}

Decentralized Finance (DeFi) changed how capital is allocated on public blockchains: lending, borrowing, and leveraged trading can run without traditional financial intermediaries \cend{1}. Yet most credit-like interactions remain highly collateralized because protocols cannot rely on off-chain identity systems or conventional credit histories \cend{2}. In a pseudonymous environment, the operational challenge is to allocate capital safely when counterparties are addresses rather than identified persons.

This research addresses a concrete research problem within that setting: whether an on-chain reputation score built from ERC-20 \texttt{approve}/\texttt{permit} allowances---implemented as EndorseRank---can reduce computational cost, represent endorsement semantics more faithfully, and yield a distinct yet credible alignment profile compared with transfer-based Adaptive Weighted PageRank (AWP) on Arbitrum One. The primary empirical comparison is EndorseRank versus AWP on an identical wallet sample, checked with contemporaneous transfer and allowance statistics and with a temporal holdout of future new approvals, under the rank-correlation validation protocol introduced with AWP\cend{3}.

In this study, on-chain reputation score computation refers to deriving a quantitative wallet score from observable on-chain behavior and relationships without off-chain identity \cend{4}. EndorseRank produces such a score from allowance-based endorsement edges. The remainder of this chapter situates the problem historically and economically, motivates allowance-based reputation, states the research problem, objectives, and questions, defines key constructs, summarizes the literature and the methodology that later chapters develop, and outlines contributions, scope, and the research roadmap.

\dissertationSubheading[sec:background]{Background: DeFi, collateralization, and pseudonymity}

Public blockchains introduced programmable settlement and open participation without a central operator \cend{5}. Smart contracts encode agreement terms as self-executing code \cend{6}, enabling financial applications that settle on-chain rather than through custodial intermediaries \cend{7}. DeFi protocols compose these contracts into markets for lending, automated market making, and perpetual futures \cend{8}. On Arbitrum One and similar Layer-2 networks, activity density is high enough that wallet-level behavioral signals---transfers, approvals, and protocol events---can be observed at scale.

Despite this transparency, DeFi credit markets remain structurally over-collateralized. Protocols typically require borrowers or traders to post collateral exceeding the value of their exposure because counterparties are pseudonymous addresses rather than identified legal persons \cend{9,10}. Pseudonymity is a design feature of public ledgers: anyone may create addresses at negligible cost, and a single economic actor may control many wallets \cend{11}. Traditional credit scoring depends on identity-linked histories that do not transfer cleanly to this setting \cend{2}. Consequently, risk management defaults to collateral buffers, liquidation engines, and protocol-level parameters rather than individualized reputation.

Collateralization protects protocols but imposes capital inefficiency. Traders and borrowers must lock assets that could otherwise be deployed elsewhere; protocols forgo under-collateralized products that traditional finance offers to borrowers with established credit. Researchers have therefore explored on-chain signals as substitutes for off-chain identity: transfer-graph centrality \cend{4}, network risk propagation \cend{12}, prospect-theory and blockchain-based credit models \cend{13}, and machine-learning pipelines on lending protocols \cend{14}. These approaches share a premise: public ledger data encode economically meaningful relationships that can inform risk assessment without deanonymizing users.

Graph ranking methods are a natural fit for this premise. PageRank and its variants propagate influence through directed edges, ranking nodes by recursive endorsement from other nodes \cend{15}. Early web-trust systems such as EigenTrust \cend{16} and TrustRank \cend{17} applied similar ideas to peer-to-peer and web spam settings. Do and Do\cend{4} apply PageRank and HodgeRank to Ethereum transfer graphs as a social-credit measure. Adaptive Weighted PageRank (AWP), proposed by Do, Do, and Nguyen\cend{3}, is a separate transfer-graph ranking method that uses value-weighted, logistically time-decayed edges and validates scores against inbound transfer count and value as domain-intuition proxies for economic importance. That line of work establishes transfer-graph reputation as a credible baseline. It does not model explicit spending authorization, a distinct on-chain act recorded by ERC-20 allowances \cend{18,19}.

\dissertationSubheading[sec:motivation]{Motivation for on-chain reputation}

Four practical pressures motivate on-chain reputation research in leveraged DeFi.
\begin{itemize}
\item Capital efficiency: protocols seek ordinal signals that could support differentiated service---for example, limit-increase cadence, fee tiers, or monitoring priority---and, only under future controlled experiments, modest collateral differentiation, while preserving on-chain settlement \cend{1}.
\item Auditability: unlike opaque off-chain credit scores, graph-based scores derived from public events can be recomputed and inspected, addressing concerns about black-box decentralized credit scoring \cend{2}.
\item Operational cost: reputation systems that must refresh frequently---for example, before each leverage adjustment---cannot afford expensive historical transfer aggregation on every update.
\item Standardization: pseudonymous DeFi lacks a shared ordinal vocabulary for wallet-level counterparty differentiation of the kind that screening ratios provide in other markets. Chapter~\ref{ch:discussion} develops industry adoption scenarios in which EndorseRank and AWP play that infrastructural role---supporting fees, utilization, and monitoring efficiency---without substituting for collateral.
\end{itemize}

Transfer-based AWP addresses the first two pressures but is costly on the third. Constructing $G_{\text{transfer}}$ requires aggregating transfer history, applying logistic time-decay, and running PageRank on a dense edge set \cend{3}. Allowance graphs offer a structurally different alternative. An ERC-20 \texttt{approve} or EIP-2612 \texttt{permit} records how much a token owner authorizes a spender to transfer \cend{18,19}. The latest allowance per owner--spender pair is a snapshot of current authorization intent: newer approvals override older ones, so temporal decay is unnecessary. Because approvals are rarer on-chain events than transfers, the resulting endorsement graph $G_{\text{approve}}$ has roughly an order of magnitude fewer edges than $G_{\text{transfer}}$ on the same wallets, and PageRank---whose per-iteration cost is $O(|E|)$---completes correspondingly faster (Table~\ref{tab:benchmark-runtime}). That gap is an expected consequence of graph sparsity under latest-allowance snapshots, not an algorithmic contribution: the solver is identical for both methods. Claim~C1 reports it as an operational property of the representation.

Beyond efficiency, allowances carry distinct semantics. A transfer may reflect routing, market making, or passive receipt; an allowance is an explicit delegation of spending authority. Treating allowances as endorsement edges aligns PageRank's citation metaphor \cend{15} with a trust-like act that protocols already rely on for token interactions. Whether that semantic difference produces rankings that differ, and whether those rankings align differently with authorization and transfer checks, is an empirical question. This research answers it on Arbitrum One using a matched sample of wallets with non-zero endorsement- and transfer-graph connectivity ($n=5{,}521$).\footnote{The matched cohort is wallets with non-zero connectivity in the endorsement and transfer subgraphs in the observation window. Runtime sensitivity to node count uses a larger expanded wallet pool ($N=99{,}080$).}

\dissertationSubheading[sec:research-problem]{Research problem}

Existing Adaptive Weighted PageRank (AWP) methods infer trust from transfer volume and logistic time-decay over transaction history \cend{3}. While economically interpretable, they (i) require expensive historical transfer aggregation, (ii) treat inbound transfer volume as a proxy for a ``valuable wallet,'' and (iii) do not model explicit spending authorization. EndorseRank instead uses the latest allowance per owner--spender pair as edge strength on a directed endorsement graph, yielding a sparser graph and a semantically distinct trust signal without temporal decay weighting.

The research problem is therefore comparative and evaluative: on an identical Arbitrum One wallet sample, does EndorseRank produce wallet orderings that are construct-distinct from AWP, align with allowance and transfer checks, predict later new endorsements in a time holdout, and do so at substantially lower computational cost? The primary comparison is EndorseRank versus AWP.

Table~\ref{tab:primary-methods-preview} summarizes the primary comparison. Table~\ref{tab:hrq-claim-map} summarizes hypotheses (H1--H3), research questions (RQ1--RQ4), and Chapter~4 empirical claims (C1--C4).

\begin{table}[htbp]
\centering
\caption{Primary comparison: EndorseRank versus AWP (matched wallet cohort, $n=5{,}521$).}
\label{tab:primary-methods-preview}
\fitwidth{%
\begin{tabular}{p{0.18\linewidth}p{0.28\linewidth}p{0.48\linewidth}}
\toprule
Method & Input graph & Role \\
\midrule
EndorseRank & Latest ERC-20 allowances ($G_{\text{approve}}$) & Primary construct: allowance delegation; target of this research \\
AWP & Time-decayed transfers ($G_{\text{transfer}}$) & Primary baseline: IEEE RIVF 2023 transfer-graph method \cend{3}; Do and Do\cend{4} is the earlier PageRank/HodgeRank study on Ethereum transfers \\
\bottomrule
\end{tabular}
}
\end{table}

To avoid ambiguity, the key concepts used throughout this research are defined as follows:

\begin{itemize}
\item \emph{DeFi (Decentralized Finance):} Financial services---including lending and borrowing---implemented via smart contracts on public blockchains, typically without a central operator controlling custody or settlement \cend{1}. DeFi protocols compose open interfaces so that token movements are observable on-chain.
\item \emph{Smart contract:} Self-executing program code deployed on a blockchain that enforces agreement terms without a trusted intermediary \cend{6}. In this research, smart contracts are the source of ERC-20 allowance logs and transfer logs.
\item \emph{Wallet / Address:} A blockchain account that holds assets and interacts with smart contracts; the unit of analysis for reputation scoring \cend{7}. A wallet is identified by its address; no off-chain identity is assumed or recovered.
\item \emph{On-chain reputation score:} A numerical indicator derived from on-chain signals that represents a wallet's relative standing for financial interactions under pseudonymity \cend{4}. Scores are ordinal for ranking purposes; absolute magnitudes are not interpreted as credit limits.
\item \emph{PageRank:} A graph ranking method that propagates influence through directed edges under a damping factor \cend{15}. EndorseRank applies PageRank to a latest-allowance endorsement graph: edge strength equals the current allowance per owner--spender pair, with no temporal decay. Adaptive Weighted PageRank (AWP) applies PageRank to transfer graphs with edge weights given by transfer value and logistic time-decay \cend{3}.
\item \emph{ERC-20 allowance / Approve / Permit:} Mechanisms recording how much a token owner permits a spender to transfer on the owner's behalf \cend{18,19}. \texttt{Approve} is a standard on-chain call; EIP-2612 \texttt{permit} authorizes the same state change via signature. EndorseRank uses the latest non-zero allowance snapshot per owner--spender pair.
\item \emph{Endorsement graph:} A directed graph in which edges represent explicit delegation of spending authority rather than mere asset transfers \cend{18}. In EndorseRank, an edge $u \to v$ means wallet $u$ currently authorizes wallet or contract $v$ to spend tokens owned by $u$; transfer-graph AWP provides the contrasting baseline \cend{3}.
\item \emph{Domain alignment:} The degree to which a reputation ranking agrees with an observable proxy ranking, assessed by Spearman's $\rho$ \cend{20} and Kendall's $\tau$ \cend{21}, in the construct-validity sense of Cronbach and Meehl\cend{22}. Alignment is construct-specific: a method may align strongly with its primary graph construct and weakly with an orthogonal check.
\end{itemize}

\dissertationSubheading[sec:supplementary-extension]{Temporal holdout}

A contemporaneous allowance Kendall $\tau$ (Table~\ref{tab:alignment-allowance}) is an intended-construct check: PageRank on $G_{\text{approve}}$ is expected to agree with spender-side approval degree because the score is a global smoothing of that degree. To test whether the score carries information beyond the same-window graph, scores are frozen at 28~February~2026 and compared with new owner--spender approval pairs that first appear in March--May~2026. The claimed holdout cohort is t1 inbound spenders ($n=1{,}335$).

\begin{table}[htbp]
\centering
\caption{Summary of hypotheses, research questions, and empirical claims.}
\label{tab:hrq-claim-map}
\fitwidth{%
\begin{tabular}{p{0.12\linewidth}p{0.42\linewidth}p{0.40\linewidth}}
\toprule
Item & Addresses & Reported in \\
\midrule
H1 / RQ1 & Do EndorseRank and AWP yield distinct wallet orderings? & \S\ref{sec:rank-divergence}; inter-method $\tau$ with 95\% CI \\
H2 / RQ2 & Which social method aligns with allowance vs.\ transfer checks? & Tables~\ref{tab:alignment-family-ci} and~\ref{tab:tau-diff}; Claims C2--C3 \\
H3 / RQ3 & Is EndorseRank faster than AWP at equal $n$? & Tables~\ref{tab:benchmark-runtime}--\ref{tab:benchmark-scaling}; Claims C1, C4 \\
RQ4 & Do t1 scores predict t2 new approvals beyond same-window allowance? & \S\ref{sec:gf-pr-ceiling}; Table~\ref{tab:holdout-spenders} \\
C1 & EndorseRank vs.\ AWP runtime and graph size & Table~\ref{tab:benchmark-runtime} \\
C2 & Intended-construct checks: allowance for ER, transfer for AWP & Tables~\ref{tab:alignment-transfer}--\ref{tab:alignment-allowance} \\
C3 & Construct-specific winners (EndorseRank vs.\ AWP) & Table~\ref{tab:alignment-family-ci} \\
C4 & Efficiency--alignment trade-off for EndorseRank & Tables~\ref{tab:benchmark-runtime}, \ref{tab:alignment-family-ci} and~\ref{tab:holdout-spenders} \\
\bottomrule
\end{tabular}
}
\end{table}

\dissertationSubheading[sec:research-objectives]{Research objectives}

The study has four primary objectives, mapped one-to-one onto RQ1--RQ4. Each objective names the measurement that operationalizes it and the criterion by which Chapter~\ref{ch:discussion} judges whether it has been met.

\begin{enumerate}
\item[O1.] Determine whether EndorseRank and AWP induce distinct wallet orderings on an identical Arbitrum One cohort (RQ1). Measurement: Kendall $\tau$ between the two score vectors on the matched cohort ($n=5{,}521$) with a 95\% bootstrap confidence interval. Criterion: the interval excludes both $0$ (unrelated orderings) and $1$ (identical orderings).
\item[O2.] Measure contemporaneous construct checks---allowance for EndorseRank and transfer for AWP---on the same cohort (RQ2). Measurement: family-level Kendall $\tau$ with 95\% intervals and paired bootstrap contrasts $\Delta\tau$ between methods within each family. Criterion: each method's $\Delta\tau$ on its own construct is positive with an interval excluding $0$.
\item[O3.] Compare PageRank computational cost of EndorseRank and AWP at equal $n$, including node-count sensitivity on the expanded pool (RQ3). Measurement: wall-clock solve time (mean and SD over five timed repeats after one warm-up), peak memory, iterations, and $|V|$, $|E|$ per configuration under one fixed protocol. Criterion: the runtime ratio is reported together with the edge ratio so that any speed advantage is attributed to sparsity.
\item[O4.] Test whether t1 scores predict t2 new owner--spender approvals on the spender cohort (RQ4). Measurement: Kendall $\tau$ between scores frozen at 28~February~2026 and new-approval counts in March--May~2026 on $n=1{,}335$ inbound spenders, with a 95\% interval, compared against AWP and against the same-window allowance degree. Criterion: EndorseRank's interval excludes $0$ and its point estimate exceeds the AWP estimate.
\end{enumerate}

\dissertationSubheading[sec:research-questions]{Research questions}

This study tests three primary hypotheses: H1---EndorseRank and AWP yield distinct wallet orderings on the same cohort; H2---the two social methods differ in construct-specific checks (allowance versus transfer); and H3---EndorseRank is computationally more efficient than AWP at equal $n$. The following research questions guide the empirical evaluation summarized in Table~\ref{tab:hrq-claim-map}.

\begin{enumerate}
\item[RQ1.] Do EndorseRank and AWP produce meaningfully different wallet orderings on an identical Arbitrum wallet sample?
\item[RQ2.] Between EndorseRank and AWP, which social method aligns more strongly with contemporaneous allowance checks and with contemporaneous transfer checks (Spearman's $\rho$, Kendall's $\tau$)?
\item[RQ3.] Is EndorseRank computationally more efficient than AWP on the same matched sample (runtime, memory, iterations)?
\item[RQ4.] After scores are frozen at t1, does EndorseRank predict t2 new owner--spender approvals on inbound spenders more strongly than AWP, and more informatively than the same-window allowance check?
\end{enumerate}

RQ1 tests whether allowance and transfer graphs induce distinct centrality structures, operationalized by inter-method Kendall $\tau$ between score vectors. RQ2 evaluates intended-construct checks on transfer and allowance local statistics; because each method's own-construct proxies are derived from the same event stream that builds its graph, RQ2 is read as a construct-validity observation, not as external validation. RQ3 compares PageRank wall-clock time, memory, and iterations on identical wallet sets, including node-count sensitivity on an expanded wallet pool. RQ4 is a time-split test of the same endorsement construct: future new approvals, not trading profit. RQ4 replaces the trading-success comparison of the research proposal, which was withdrawn because the outcome-native reference it relied on was circular with the proxies (Section~\ref{sub:gf-pr}).

\dissertationSubheading[sec:literature-summary]{Summary of the literature review}

Chapter~\ref{ch:literature} reviews four bodies of work; the following summary states what each contributes to the design and where it falls short. First, graph ranking. PageRank \cend{15} ranks nodes by recursively propagated endorsement and is computed by power iteration whose per-iteration cost is linear in the number of edges \cend{23}; EigenTrust \cend{16} and TrustRank \cend{17} adapt the idea to peer-to-peer trust and web spam with normalisation and seed sets that limit inflation. These methods supply the solver and the complexity argument used here, but their edge semantics are fixed to hyperlinks, ratings, or seeds. Second, on-chain reputation. Do and Do\cend{4} apply PageRank and HodgeRank to Ethereum transaction graphs as a social-credit measure; Do, Do, and Nguyen\cend{3} propose AWP, which weights transfer edges by value and a logistic decay of age and validates the ranking against inbound transfer count and value; Lin et al.\cend{12} propagate account risk through the transaction network; Wu et al.\cend{24} show that graph embeddings of transaction data outperform per-account features for phishing detection. Surveys of blockchain reputation systems \cend{25,26} confirm that explicit ratings and transfer volume are the dominant trust signals. None of these works uses the ERC-20 allowance---the on-chain act that records a deliberate delegation of spending authority \cend{18}---as the graph edge, and the AWP validation compares scores with proxies computed from the same transfer data that builds the graph. Third, DeFi risk and measurement. Systematizations of DeFi \cend{8} and of DeFi attacks \cend{27} locate behavioral counterparty risk as one layer among smart-contract, oracle, market, and governance risk; wash trading \cend{28,29}, extractable value \cend{30}, scam tokens \cend{31}, and airdrop farming \cend{32} show how much transfer activity is artificial, which is the empirical reason to look for a signal other than payment flow. Fourth, measurement theory and inference. Construct validity \cend{22,33} distinguishes what a score is designed to measure from external criteria, and bootstrap resampling \cend{34,35} provides distribution-free intervals for rank statistics. The gap that emerges (Section~\ref{sec:research-gap}) is a same-cohort comparison of an allowance-edge ranking against the transfer-edge baseline, with confidence intervals, paired difference tests, a fixed timing protocol that reports edge counts, and a temporal holdout that separates same-window agreement from prediction.

\dissertationSubheading[sec:methodology-summary]{Summary of the research methodology}

Chapter~\ref{ch:methodology} adopts a quantitative, comparative design on public ledger data. ERC-20 \texttt{Approval} and \texttt{Transfer} logs for Arbitrum One over 1~December~2025--31~May~2026 are extracted from Google BigQuery; the matched wallet cohort ($n=5{,}521$) comprises addresses with non-zero connectivity in both the latest-allowance graph $G_{\text{approve}}$ and the time-decayed transfer graph $G_{\text{transfer}}$. Both graphs are scored with one shared power-iteration PageRank implementation ($d=0.85$, $\varepsilon=10^{-8}$, at most 300 iterations); EndorseRank uses summed latest allowances as edge weights, AWP uses value-weighted, logistically decayed transfers exactly as specified in the baseline paper\cend{3}. Rank agreement is measured with Spearman's $\rho$ and Kendall's $\tau$ against three contemporaneous proxy families (transfer, allowance, Sybil-adjusted stability) and, for RQ4, against new owner--spender approvals in a three-month holdout after scores are frozen. Every coefficient carries a 95\% percentile bootstrap interval from 400 paired wallet resamples (seed 42), and pre-specified method contrasts are reported as paired $\Delta\tau$ intervals. Runtime, peak memory, iterations, and $|V|$, $|E|$ are recorded under one timing protocol (one untimed warm-up, five timed repeats, per-session memoization so that identical graphs yield identical numbers wherever they are tabulated). Sensitivity analyses vary damping, restrict to top tokens, and change the wallet-sample definition. The pipeline, configuration, and machine-readable results summary from which every table is generated are versioned in a public repository (Section~\ref{sec:reproducibility}); the study uses public data only and involves no human participants (Section~\ref{sec:ethics}).

\dissertationSubheading[sec:contributions]{Contributions}

This research makes three contributions to on-chain reputation research and DeFi risk assessment. EndorseRank applies classical PageRank \cend{15} unchanged. The contribution is the combination of a new edge construct, a same-cohort comparison with inference, and an out-of-window test.
\begin{itemize}
\item Latest-allowance edge construction (EndorseRank): an endorsement graph in which the edge $u\to v$ carries the current ERC-20 \texttt{approve}/\texttt{permit} allowance from owner $u$ to spender $v$. Newer approvals override older ones, so the construct needs no temporal decay, and it records explicit delegation of spending authority rather than payment flow. No reviewed reputation method uses this edge (Section~\ref{sec:research-gap}).
\item Same-cohort contemporaneous comparison with inference: EndorseRank and AWP are computed on one matched wallet cohort and compared on identical transfer, allowance, and Sybil-stability proxies under the AWP validation protocol\cend{3}, with 95\% bootstrap intervals on every coefficient and paired $\Delta\tau$ contrasts \cend{34}. The design states that each method's own-construct agreement is partly mechanical, and it reports edge counts alongside runtimes so that the efficiency result (Claim~C1) is read as an expected consequence of sparsity.
\item Temporal holdout of future new approvals: scores frozen at t1 are tested against new owner--spender approvals in t2 on the spender cohort, separating same-window agreement from prediction. Together with the open, versioned evaluation pipeline (Section~\ref{sec:reproducibility}; Appendix~\ref{ch:appendix-eval}), this provides the evidence base for Claims~C1--C4.
\end{itemize}

Headline numerical results for inter-method $\tau$, family-level alignment, and runtime are reported once in Chapter~\ref{ch:results} (Tables~\ref{tab:benchmark-runtime} and~\ref{tab:alignment-family-ci}) and are not restated here. Those results describe an efficiency--alignment trade-off: neither social method dominates both constructs.

\dissertationSubheading[sec:scope-limitations]{Scope and limitations}

This research focuses on Arbitrum One over the observation window 2025-12-01 to 2026-05-31. The primary empirical sample is a matched wallet cohort of $n=5{,}521$ addresses with non-zero endorsement- and transfer-graph connectivity; edge counts for both graphs are reported in Table~\ref{tab:benchmark-runtime}. Runtime sensitivity to node count is reported on an expanded wallet pool ($N=99{,}080$). Sensitivity analyses examine damping, top-token subgraphs, and the wallet-sample definition.

Unsecured lending default labels are out of scope. Constructs and ethics are defined in Chapters~\ref{ch:introduction} and~\ref{ch:methodology}. The holdout claim uses the spender cohort ($n=1{,}335$). Extended baselines archived for future comparison are excluded from the main narrative (Appendix~\ref{ch:appendix-eval}).

Generalization beyond Arbitrum One and the six-month window is not claimed. Token-decimal normalization, price oracles, and permit-versus-approve coverage may affect edge weights. Adversarial manipulation of allowances or transfers remains a threat discussed in Chapter~\ref{ch:discussion}; Sybil-adjusted proxies partially address but do not eliminate that risk \cend{11}.

\dissertationSubheading[sec:roadmap]{Research roadmap}

This research proceeds as follows.

Chapter~\ref{ch:literature} reviews graph-based on-chain reputation and ERC-20 allowance semantics. It develops the theoretical framework---hyperlink propagation, domain alignment, and computational complexity---that motivates the primary EndorseRank versus AWP comparison.

Chapter~\ref{ch:methodology} specifies data sources, sampling, graph construction for EndorseRank and AWP, contemporaneous construct checks, the temporal holdout, computational benchmarks, robustness checks, reproducibility practices, and ethics. It defines the matched wallet cohort, the spender holdout cohort, the expanded wallet pool used for scaling, and the operational procedures that produce the Chapter~\ref{ch:results} tables.

Chapter~\ref{ch:results} reports empirical results: cohort characteristics, runtime and scaling benchmarks (Claim~C1), robustness checks, allowance and transfer alignment (Claims~C2--C3), rank divergence (H1/RQ1), and the temporal holdout (RQ4). Primary claims C1--C4 are stated explicitly.

Chapter~\ref{ch:discussion} interprets findings by research question, develops theoretical and practical implications, formulates a Sybil attack cost model on allowance graphs, discusses regulatory and ethical considerations, and evaluates threats to validity and limitations.

Chapter~\ref{ch:conclusion} synthesizes the contribution to knowledge, answers RQ1--RQ4 in turn, and outlines future work on cross-protocol validation, adversarial robustness, hybrid methods, and live oracle integration.

Chapter~\ref{ch:literature} starts from the blockchain and DeFi primitives that the allowance construct uses and closes with the research gap this study is designed to close.
